\documentclass[12pt]{article}
\usepackage[utf8]{inputenc}
\usepackage[T1]{fontenc}
\usepackage{hyperref}
\title{Bigambal Language Resources for K-12 Curriculum Development}
\author{Compiled by Claude Sonnet 5.5 (claude-sonnet-5-5) for Andr\'as Kornai}
\date{}
\begin{document}
\maketitle

\section{Introduction}
Bigambal (also spelled Bigambul, Pikambul; ISO 639-3 xbe) is an extinct Pama-Nyungan language of the Goondiwindi region on the Queensland--New South Wales border. Its classification is unresolved: Dixon (2002) links it with Bundjalung, while O'Grady, Voegelin and Voegelin class it as Wiradjuric. AIATSIS lists Gambuwal and Kwiambal as dialects. The search found very little material, and no curriculum resources.

\section{Primary Educational Resources}
No primers, textbooks or school programmes for Bigambal were found in search results.

\section{Linguistic and Reference Materials}
\begin{itemize}
\item AIATSIS AUSTLANG entry D34 (Bigambul): language region, alternative names and classification. \url{https://collection.aiatsis.gov.au/austlang/language/d34}
\item Wikipedia, ``Bigambal language'': a short overview with a few recorded words, such as \emph{bamburr} (kangaroo) and \emph{warril} (river). \url{https://en.wikipedia.org/wiki/Bigambal_language}
\item Wikipedia, ``Bigambul language''. \url{https://en.wikipedia.org/wiki/Bigambul_language}
\end{itemize}
No published Bigambal dictionary or grammar was found.

\section{Religious and Community Resources}
A 2020 volunteer report from Augustinian Justice and Peace says a group has worked with the Bigambul Aboriginal Corporation of Goondiwindi to reconstruct the language. It says several hundred words were recovered from early pastoralists', explorers' and pastors' lists, reorganised by linguist Jeremy Steele, and that recordings were made so the community can compare pronunciations. \url{https://osa.org.au/?p=4927}

\section{Digital and Audio Resources}
Recordings are mentioned in the report above, but no publicly accessible audio, video or software was found.

\section{Curriculum Development Context}
Material is limited to archival word lists and a community reclamation effort reported in 2020. Any curriculum would need to be developed with the Bigambul community and the Bigambul Aboriginal Corporation. Because Gamilaraay may have been spoken in the southern part of the territory, the sources should be checked carefully. I found no status updates after 2020.
\end{document}